% paper15-theorem-ecologies.tex — Theorem Ecologies: Evolving Populations of Verification Lemmas
% Paper 15 in the Judgment Geometry series.
% Compile: pdflatex paper15-theorem-ecologies && pdflatex paper15-theorem-ecologies
\documentclass[11pt]{article}
% jugeo-common.tex — shared preamble for all Judgment Geometry papers
% Include via: % jugeo-common.tex — shared preamble for all Judgment Geometry papers
% Include via: % jugeo-common.tex — shared preamble for all Judgment Geometry papers
% Include via: \input{jugeo-common}

% ── Packages ────────────────────────────────────────────────────────────
\usepackage{amsmath,amssymb,amsthm,mathtools}
\usepackage{tikz-cd}
\usepackage{listings}
\usepackage{xcolor}
\usepackage{xspace}
\usepackage{booktabs}
\usepackage{multirow}
\usepackage{enumitem}
\usepackage[expansion=false]{microtype}
\usepackage{hyperref}
\usepackage{cleveref}
% \usepackage{thmtools}  % disabled: not available in this TeX installation
\usepackage{float}
\usepackage{caption}
\usepackage{subcaption}
\usepackage{tabularx}
\usepackage{stmaryrd}       % \llbracket, \rrbracket
\usepackage{bbm}            % \mathbbm{1}

% ── Page geometry ───────────────────────────────────────────────────────
\usepackage[margin=1in]{geometry}

% ── Theorem environments ────────────────────────────────────────────────
\theoremstyle{plain}
\newtheorem{theorem}{Theorem}[section]
\newtheorem{lemma}[theorem]{Lemma}
\newtheorem{proposition}[theorem]{Proposition}
\newtheorem{corollary}[theorem]{Corollary}

\theoremstyle{definition}
\newtheorem{definition}[theorem]{Definition}
\newtheorem{example}[theorem]{Example}
\newtheorem{construction}[theorem]{Construction}

\theoremstyle{remark}
\newtheorem{remark}[theorem]{Remark}
\newtheorem{notation}[theorem]{Notation}

% ── Python listings style ──────────────────────────────────────────────
\definecolor{jugeoBlue}{HTML}{2563EB}
\definecolor{jugeoGreen}{HTML}{059669}
\definecolor{jugeoRed}{HTML}{DC2626}
\definecolor{jugeoPurple}{HTML}{7C3AED}
\definecolor{jugeoGray}{HTML}{6B7280}
\definecolor{jugeoBg}{HTML}{F8FAFC}

\lstdefinestyle{jugeo-python}{
  language=Python,
  basicstyle=\ttfamily\small,
  keywordstyle=\color{jugeoBlue}\bfseries,
  stringstyle=\color{jugeoGreen},
  commentstyle=\color{jugeoGray}\itshape,
  numberstyle=\tiny\color{jugeoGray},
  backgroundcolor=\color{jugeoBg},
  numbers=left,
  numbersep=6pt,
  frame=single,
  framerule=0.4pt,
  rulecolor=\color{jugeoGray!40},
  breaklines=true,
  breakatwhitespace=true,
  showstringspaces=false,
  tabsize=4,
  xleftmargin=1.5em,
  framexleftmargin=1.5em,
  morekeywords={dataclass,frozen,field,Enum,IntEnum,auto,Any,
                Mapping,Sequence,Optional,match,case},
  literate={→}{{$\to$}}1 {←}{{$\leftarrow$}}1
           {≤}{{$\leq$}}1 {≥}{{$\geq$}}1 {≠}{{$\neq$}}1
           {∈}{{$\in$}}1 {∀}{{$\forall$}}1 {∃}{{$\exists$}}1
           {⊕}{{$\oplus$}}1 {⊖}{{$\ominus$}}1,
  escapeinside={(*@}{@*)},
}
\lstset{style=jugeo-python}

\lstdefinestyle{jugeo-output}{
  basicstyle=\ttfamily\small,
  backgroundcolor=\color{jugeoBg},
  frame=single,
  framerule=0.4pt,
  rulecolor=\color{jugeoGray!40},
  breaklines=true,
  showstringspaces=false,
  xleftmargin=1.5em,
  framexleftmargin=0.5em,
  numbers=none,
}

% ── JuGeo-specific macros ──────────────────────────────────────────────

% Judgment 8-tuple
\newcommand{\J}{\mathcal{J}}
\newcommand{\Jud}[8]{(#1,\,#2,\,#3,\,#4,\,#5,\,#6,\,#7,\,#8)}
\newcommand{\JudShort}[1]{\J_{#1}}

% Components
\newcommand{\coord}{\mathsf{c}}            % coordinate
\newcommand{\prop}{\varphi}                 % proposition
\newcommand{\carrier}{\mathcal{A}}          % carrier type
\newcommand{\evid}{\mathcal{E}}             % evidence bundle
\newcommand{\oblig}{\mathcal{O}}            % obligations
\newcommand{\obstr}{\mathcal{B}}            % obstructions
\newcommand{\trust}{\mathcal{T}}            % trust annotation
\newcommand{\prov}{\Pi}                     % provenance

% Trust algebra
\newcommand{\Talg}{\mathfrak{T}}            % trust ordered algebra
\newcommand{\Eadm}{\mathcal{E}_{\mathrm{adm}}}
\newcommand{\tleq}{\preceq}                 % trust ordering
\newcommand{\tjoin}{\oplus}                 % conservative join
\newcommand{\tatten}{\ominus}               % attenuation
\newcommand{\tpromote}{\uparrow_\pi}        % promotion
\newcommand{\tdemote}{\downarrow_\chi}       % demotion

% Trust tiers
\newcommand{\Tcontra}{\ensuremath{\mathsf{CONTRADICTED}}}
\newcommand{\Tunver}{\ensuremath{\mathsf{UNVERIFIED}}}
\newcommand{\Tcopilot}{\ensuremath{\mathsf{COPILOT}}}
\newcommand{\Toracle}{\ensuremath{\mathsf{ORACLE}}}
\newcommand{\Truntime}{\ensuremath{\mathsf{RUNTIME}}}
\newcommand{\Tsolver}{\ensuremath{\mathsf{SOLVER}}}
\newcommand{\Tproof}{\ensuremath{\mathsf{PROOF}}}

% Site and geometry
\newcommand{\Site}{\mathcal{S}}             % semantic site
\newcommand{\Cov}{\mathrm{Cov}}            % covering family
\newcommand{\Ob}{\mathrm{Ob}}              % objects
\newcommand{\Mor}{\mathrm{Mor}}            % morphisms
\newcommand{\res}{\mathrm{res}}            % restriction
\newcommand{\Cech}{\check{C}}              % Čech complex
\newcommand{\Hcoh}[1]{H^{#1}}             % cohomology group

% Descent
\newcommand{\descent}{\mathrm{desc}}
\newcommand{\glue}{\mathrm{glue}}
\newcommand{\overlap}[2]{#1 \cap #2}

% Semantic moves
\newcommand{\VERIFY}{\textsc{Verify}}
\newcommand{\CONSTRUCT}{\textsc{Construct}}
\newcommand{\NORMALIZE}{\textsc{Normalize}}
\newcommand{\GLUE}{\textsc{Glue}}
\newcommand{\RESTRICT}{\textsc{Restrict}}
\newcommand{\TRANSPORT}{\textsc{Transport}}
\newcommand{\REPAIR}{\textsc{Repair}}
\newcommand{\PROMOTE}{\textsc{Promote}}
\newcommand{\CHALLENGE}{\textsc{Challenge}}
\newcommand{\ESCALATE}{\textsc{Escalate}}

% Fragments
\newcommand{\QFLIA}{\texttt{QF\_LIA}}
\newcommand{\QFLRA}{\texttt{QF\_LRA}}
\newcommand{\QFBV}{\texttt{QF\_BV}}
\newcommand{\QFUF}{\texttt{QF\_UF}}

% Misc
\newcommand{\Python}{\textsc{Python}}
\newcommand{\JuGeo}{\textsc{JuGeo}}
\newcommand{\LEAN}{\textsc{Lean}\,4}
\newcommand{\Fstar}{F$^{\star}$}
\newcommand{\Coq}{\textsc{Coq}}
\newcommand{\Dafny}{\textsc{Dafny}}

% Common shortcuts
\newcommand{\ie}{i.e.\@\xspace}
\newcommand{\eg}{e.g.\@\xspace}
\newcommand{\cf}{cf.\@\xspace}
\newcommand{\wrt}{w.r.t.\@\xspace}

% Evaluation shortcuts
\newcommand{\numcases}{300}
\newcommand{\numspec}{100}
\newcommand{\numeq}{100}
\newcommand{\numbug}{100}
\newcommand{\medtime}{6.5\,ms}
\newcommand{\totaltime}{1.949\,s}


% ── Packages ────────────────────────────────────────────────────────────
\usepackage{amsmath,amssymb,amsthm,mathtools}
\usepackage{tikz-cd}
\usepackage{listings}
\usepackage{xcolor}
\usepackage{xspace}
\usepackage{booktabs}
\usepackage{multirow}
\usepackage{enumitem}
\usepackage[expansion=false]{microtype}
\usepackage{hyperref}
\usepackage{cleveref}
% \usepackage{thmtools}  % disabled: not available in this TeX installation
\usepackage{float}
\usepackage{caption}
\usepackage{subcaption}
\usepackage{tabularx}
\usepackage{stmaryrd}       % \llbracket, \rrbracket
\usepackage{bbm}            % \mathbbm{1}

% ── Page geometry ───────────────────────────────────────────────────────
\usepackage[margin=1in]{geometry}

% ── Theorem environments ────────────────────────────────────────────────
\theoremstyle{plain}
\newtheorem{theorem}{Theorem}[section]
\newtheorem{lemma}[theorem]{Lemma}
\newtheorem{proposition}[theorem]{Proposition}
\newtheorem{corollary}[theorem]{Corollary}

\theoremstyle{definition}
\newtheorem{definition}[theorem]{Definition}
\newtheorem{example}[theorem]{Example}
\newtheorem{construction}[theorem]{Construction}

\theoremstyle{remark}
\newtheorem{remark}[theorem]{Remark}
\newtheorem{notation}[theorem]{Notation}

% ── Python listings style ──────────────────────────────────────────────
\definecolor{jugeoBlue}{HTML}{2563EB}
\definecolor{jugeoGreen}{HTML}{059669}
\definecolor{jugeoRed}{HTML}{DC2626}
\definecolor{jugeoPurple}{HTML}{7C3AED}
\definecolor{jugeoGray}{HTML}{6B7280}
\definecolor{jugeoBg}{HTML}{F8FAFC}

\lstdefinestyle{jugeo-python}{
  language=Python,
  basicstyle=\ttfamily\small,
  keywordstyle=\color{jugeoBlue}\bfseries,
  stringstyle=\color{jugeoGreen},
  commentstyle=\color{jugeoGray}\itshape,
  numberstyle=\tiny\color{jugeoGray},
  backgroundcolor=\color{jugeoBg},
  numbers=left,
  numbersep=6pt,
  frame=single,
  framerule=0.4pt,
  rulecolor=\color{jugeoGray!40},
  breaklines=true,
  breakatwhitespace=true,
  showstringspaces=false,
  tabsize=4,
  xleftmargin=1.5em,
  framexleftmargin=1.5em,
  morekeywords={dataclass,frozen,field,Enum,IntEnum,auto,Any,
                Mapping,Sequence,Optional,match,case},
  literate={→}{{$\to$}}1 {←}{{$\leftarrow$}}1
           {≤}{{$\leq$}}1 {≥}{{$\geq$}}1 {≠}{{$\neq$}}1
           {∈}{{$\in$}}1 {∀}{{$\forall$}}1 {∃}{{$\exists$}}1
           {⊕}{{$\oplus$}}1 {⊖}{{$\ominus$}}1,
  escapeinside={(*@}{@*)},
}
\lstset{style=jugeo-python}

\lstdefinestyle{jugeo-output}{
  basicstyle=\ttfamily\small,
  backgroundcolor=\color{jugeoBg},
  frame=single,
  framerule=0.4pt,
  rulecolor=\color{jugeoGray!40},
  breaklines=true,
  showstringspaces=false,
  xleftmargin=1.5em,
  framexleftmargin=0.5em,
  numbers=none,
}

% ── JuGeo-specific macros ──────────────────────────────────────────────

% Judgment 8-tuple
\newcommand{\J}{\mathcal{J}}
\newcommand{\Jud}[8]{(#1,\,#2,\,#3,\,#4,\,#5,\,#6,\,#7,\,#8)}
\newcommand{\JudShort}[1]{\J_{#1}}

% Components
\newcommand{\coord}{\mathsf{c}}            % coordinate
\newcommand{\prop}{\varphi}                 % proposition
\newcommand{\carrier}{\mathcal{A}}          % carrier type
\newcommand{\evid}{\mathcal{E}}             % evidence bundle
\newcommand{\oblig}{\mathcal{O}}            % obligations
\newcommand{\obstr}{\mathcal{B}}            % obstructions
\newcommand{\trust}{\mathcal{T}}            % trust annotation
\newcommand{\prov}{\Pi}                     % provenance

% Trust algebra
\newcommand{\Talg}{\mathfrak{T}}            % trust ordered algebra
\newcommand{\Eadm}{\mathcal{E}_{\mathrm{adm}}}
\newcommand{\tleq}{\preceq}                 % trust ordering
\newcommand{\tjoin}{\oplus}                 % conservative join
\newcommand{\tatten}{\ominus}               % attenuation
\newcommand{\tpromote}{\uparrow_\pi}        % promotion
\newcommand{\tdemote}{\downarrow_\chi}       % demotion

% Trust tiers
\newcommand{\Tcontra}{\ensuremath{\mathsf{CONTRADICTED}}}
\newcommand{\Tunver}{\ensuremath{\mathsf{UNVERIFIED}}}
\newcommand{\Tcopilot}{\ensuremath{\mathsf{COPILOT}}}
\newcommand{\Toracle}{\ensuremath{\mathsf{ORACLE}}}
\newcommand{\Truntime}{\ensuremath{\mathsf{RUNTIME}}}
\newcommand{\Tsolver}{\ensuremath{\mathsf{SOLVER}}}
\newcommand{\Tproof}{\ensuremath{\mathsf{PROOF}}}

% Site and geometry
\newcommand{\Site}{\mathcal{S}}             % semantic site
\newcommand{\Cov}{\mathrm{Cov}}            % covering family
\newcommand{\Ob}{\mathrm{Ob}}              % objects
\newcommand{\Mor}{\mathrm{Mor}}            % morphisms
\newcommand{\res}{\mathrm{res}}            % restriction
\newcommand{\Cech}{\check{C}}              % Čech complex
\newcommand{\Hcoh}[1]{H^{#1}}             % cohomology group

% Descent
\newcommand{\descent}{\mathrm{desc}}
\newcommand{\glue}{\mathrm{glue}}
\newcommand{\overlap}[2]{#1 \cap #2}

% Semantic moves
\newcommand{\VERIFY}{\textsc{Verify}}
\newcommand{\CONSTRUCT}{\textsc{Construct}}
\newcommand{\NORMALIZE}{\textsc{Normalize}}
\newcommand{\GLUE}{\textsc{Glue}}
\newcommand{\RESTRICT}{\textsc{Restrict}}
\newcommand{\TRANSPORT}{\textsc{Transport}}
\newcommand{\REPAIR}{\textsc{Repair}}
\newcommand{\PROMOTE}{\textsc{Promote}}
\newcommand{\CHALLENGE}{\textsc{Challenge}}
\newcommand{\ESCALATE}{\textsc{Escalate}}

% Fragments
\newcommand{\QFLIA}{\texttt{QF\_LIA}}
\newcommand{\QFLRA}{\texttt{QF\_LRA}}
\newcommand{\QFBV}{\texttt{QF\_BV}}
\newcommand{\QFUF}{\texttt{QF\_UF}}

% Misc
\newcommand{\Python}{\textsc{Python}}
\newcommand{\JuGeo}{\textsc{JuGeo}}
\newcommand{\LEAN}{\textsc{Lean}\,4}
\newcommand{\Fstar}{F$^{\star}$}
\newcommand{\Coq}{\textsc{Coq}}
\newcommand{\Dafny}{\textsc{Dafny}}

% Common shortcuts
\newcommand{\ie}{i.e.\@\xspace}
\newcommand{\eg}{e.g.\@\xspace}
\newcommand{\cf}{cf.\@\xspace}
\newcommand{\wrt}{w.r.t.\@\xspace}

% Evaluation shortcuts
\newcommand{\numcases}{300}
\newcommand{\numspec}{100}
\newcommand{\numeq}{100}
\newcommand{\numbug}{100}
\newcommand{\medtime}{6.5\,ms}
\newcommand{\totaltime}{1.949\,s}


% ── Packages ────────────────────────────────────────────────────────────
\usepackage{amsmath,amssymb,amsthm,mathtools}
\usepackage{tikz-cd}
\usepackage{listings}
\usepackage{xcolor}
\usepackage{xspace}
\usepackage{booktabs}
\usepackage{multirow}
\usepackage{enumitem}
\usepackage[expansion=false]{microtype}
\usepackage{hyperref}
\usepackage{cleveref}
% \usepackage{thmtools}  % disabled: not available in this TeX installation
\usepackage{float}
\usepackage{caption}
\usepackage{subcaption}
\usepackage{tabularx}
\usepackage{stmaryrd}       % \llbracket, \rrbracket
\usepackage{bbm}            % \mathbbm{1}

% ── Page geometry ───────────────────────────────────────────────────────
\usepackage[margin=1in]{geometry}

% ── Theorem environments ────────────────────────────────────────────────
\theoremstyle{plain}
\newtheorem{theorem}{Theorem}[section]
\newtheorem{lemma}[theorem]{Lemma}
\newtheorem{proposition}[theorem]{Proposition}
\newtheorem{corollary}[theorem]{Corollary}

\theoremstyle{definition}
\newtheorem{definition}[theorem]{Definition}
\newtheorem{example}[theorem]{Example}
\newtheorem{construction}[theorem]{Construction}

\theoremstyle{remark}
\newtheorem{remark}[theorem]{Remark}
\newtheorem{notation}[theorem]{Notation}

% ── Python listings style ──────────────────────────────────────────────
\definecolor{jugeoBlue}{HTML}{2563EB}
\definecolor{jugeoGreen}{HTML}{059669}
\definecolor{jugeoRed}{HTML}{DC2626}
\definecolor{jugeoPurple}{HTML}{7C3AED}
\definecolor{jugeoGray}{HTML}{6B7280}
\definecolor{jugeoBg}{HTML}{F8FAFC}

\lstdefinestyle{jugeo-python}{
  language=Python,
  basicstyle=\ttfamily\small,
  keywordstyle=\color{jugeoBlue}\bfseries,
  stringstyle=\color{jugeoGreen},
  commentstyle=\color{jugeoGray}\itshape,
  numberstyle=\tiny\color{jugeoGray},
  backgroundcolor=\color{jugeoBg},
  numbers=left,
  numbersep=6pt,
  frame=single,
  framerule=0.4pt,
  rulecolor=\color{jugeoGray!40},
  breaklines=true,
  breakatwhitespace=true,
  showstringspaces=false,
  tabsize=4,
  xleftmargin=1.5em,
  framexleftmargin=1.5em,
  morekeywords={dataclass,frozen,field,Enum,IntEnum,auto,Any,
                Mapping,Sequence,Optional,match,case},
  literate={→}{{$\to$}}1 {←}{{$\leftarrow$}}1
           {≤}{{$\leq$}}1 {≥}{{$\geq$}}1 {≠}{{$\neq$}}1
           {∈}{{$\in$}}1 {∀}{{$\forall$}}1 {∃}{{$\exists$}}1
           {⊕}{{$\oplus$}}1 {⊖}{{$\ominus$}}1,
  escapeinside={(*@}{@*)},
}
\lstset{style=jugeo-python}

\lstdefinestyle{jugeo-output}{
  basicstyle=\ttfamily\small,
  backgroundcolor=\color{jugeoBg},
  frame=single,
  framerule=0.4pt,
  rulecolor=\color{jugeoGray!40},
  breaklines=true,
  showstringspaces=false,
  xleftmargin=1.5em,
  framexleftmargin=0.5em,
  numbers=none,
}

% ── JuGeo-specific macros ──────────────────────────────────────────────

% Judgment 8-tuple
\newcommand{\J}{\mathcal{J}}
\newcommand{\Jud}[8]{(#1,\,#2,\,#3,\,#4,\,#5,\,#6,\,#7,\,#8)}
\newcommand{\JudShort}[1]{\J_{#1}}

% Components
\newcommand{\coord}{\mathsf{c}}            % coordinate
\newcommand{\prop}{\varphi}                 % proposition
\newcommand{\carrier}{\mathcal{A}}          % carrier type
\newcommand{\evid}{\mathcal{E}}             % evidence bundle
\newcommand{\oblig}{\mathcal{O}}            % obligations
\newcommand{\obstr}{\mathcal{B}}            % obstructions
\newcommand{\trust}{\mathcal{T}}            % trust annotation
\newcommand{\prov}{\Pi}                     % provenance

% Trust algebra
\newcommand{\Talg}{\mathfrak{T}}            % trust ordered algebra
\newcommand{\Eadm}{\mathcal{E}_{\mathrm{adm}}}
\newcommand{\tleq}{\preceq}                 % trust ordering
\newcommand{\tjoin}{\oplus}                 % conservative join
\newcommand{\tatten}{\ominus}               % attenuation
\newcommand{\tpromote}{\uparrow_\pi}        % promotion
\newcommand{\tdemote}{\downarrow_\chi}       % demotion

% Trust tiers
\newcommand{\Tcontra}{\ensuremath{\mathsf{CONTRADICTED}}}
\newcommand{\Tunver}{\ensuremath{\mathsf{UNVERIFIED}}}
\newcommand{\Tcopilot}{\ensuremath{\mathsf{COPILOT}}}
\newcommand{\Toracle}{\ensuremath{\mathsf{ORACLE}}}
\newcommand{\Truntime}{\ensuremath{\mathsf{RUNTIME}}}
\newcommand{\Tsolver}{\ensuremath{\mathsf{SOLVER}}}
\newcommand{\Tproof}{\ensuremath{\mathsf{PROOF}}}

% Site and geometry
\newcommand{\Site}{\mathcal{S}}             % semantic site
\newcommand{\Cov}{\mathrm{Cov}}            % covering family
\newcommand{\Ob}{\mathrm{Ob}}              % objects
\newcommand{\Mor}{\mathrm{Mor}}            % morphisms
\newcommand{\res}{\mathrm{res}}            % restriction
\newcommand{\Cech}{\check{C}}              % Čech complex
\newcommand{\Hcoh}[1]{H^{#1}}             % cohomology group

% Descent
\newcommand{\descent}{\mathrm{desc}}
\newcommand{\glue}{\mathrm{glue}}
\newcommand{\overlap}[2]{#1 \cap #2}

% Semantic moves
\newcommand{\VERIFY}{\textsc{Verify}}
\newcommand{\CONSTRUCT}{\textsc{Construct}}
\newcommand{\NORMALIZE}{\textsc{Normalize}}
\newcommand{\GLUE}{\textsc{Glue}}
\newcommand{\RESTRICT}{\textsc{Restrict}}
\newcommand{\TRANSPORT}{\textsc{Transport}}
\newcommand{\REPAIR}{\textsc{Repair}}
\newcommand{\PROMOTE}{\textsc{Promote}}
\newcommand{\CHALLENGE}{\textsc{Challenge}}
\newcommand{\ESCALATE}{\textsc{Escalate}}

% Fragments
\newcommand{\QFLIA}{\texttt{QF\_LIA}}
\newcommand{\QFLRA}{\texttt{QF\_LRA}}
\newcommand{\QFBV}{\texttt{QF\_BV}}
\newcommand{\QFUF}{\texttt{QF\_UF}}

% Misc
\newcommand{\Python}{\textsc{Python}}
\newcommand{\JuGeo}{\textsc{JuGeo}}
\newcommand{\LEAN}{\textsc{Lean}\,4}
\newcommand{\Fstar}{F$^{\star}$}
\newcommand{\Coq}{\textsc{Coq}}
\newcommand{\Dafny}{\textsc{Dafny}}

% Common shortcuts
\newcommand{\ie}{i.e.\@\xspace}
\newcommand{\eg}{e.g.\@\xspace}
\newcommand{\cf}{cf.\@\xspace}
\newcommand{\wrt}{w.r.t.\@\xspace}

% Evaluation shortcuts
\newcommand{\numcases}{300}
\newcommand{\numspec}{100}
\newcommand{\numeq}{100}
\newcommand{\numbug}{100}
\newcommand{\medtime}{6.5\,ms}
\newcommand{\totaltime}{1.949\,s}

% experiment-data.tex — AUTO-GENERATED from real jugeo CLI runs
% DO NOT EDIT — regenerate with: python3 experiments/run_universal_metrics.py
% Generated from 30 programs across 6 domains

\newcommand{\expTotalPrograms}{30}
\newcommand{\expVerified}{30}
\newcommand{\expAccuracy}{100.0\%}
\newcommand{\expCoordsMin}{2}
\newcommand{\expCoordsMax}{10}
\newcommand{\expCoordsMean}{5.2}
\newcommand{\expCoordsMedian}{5.5}
\newcommand{\expPropsMin}{4}
\newcommand{\expPropsMax}{20}
\newcommand{\expPropsMean}{10.4}
\newcommand{\expPropsSum}{311}
\newcommand{\expPropsOkSum}{310}
\newcommand{\expTimeMin}{3.506\,s}
\newcommand{\expTimeMax}{6.714\,s}
\newcommand{\expTimeMean}{4.64\,s}
\newcommand{\expTimeMedian}{4.508\,s}
\newcommand{\expTimeTotal}{139.204\,s}
\newcommand{\expObstructionTotal}{0}
\newcommand{\expHOneZero}{30}
\newcommand{\expDomainCount}{6}
\newcommand{\expTrustCOPILOTSUGGESTED}{30}
\newcommand{\expDomainDsCount}{5}
\newcommand{\expDomainDsVerified}{5}
\newcommand{\expDomainDsAvgCoords}{7.6}
\newcommand{\expDomainDsAvgProps}{15.2}
\newcommand{\expDomainMathCount}{5}
\newcommand{\expDomainMathVerified}{5}
\newcommand{\expDomainMathAvgCoords}{4.8}
\newcommand{\expDomainMathAvgProps}{9.4}
\newcommand{\expDomainSmCount}{5}
\newcommand{\expDomainSmVerified}{5}
\newcommand{\expDomainSmAvgCoords}{7.6}
\newcommand{\expDomainSmAvgProps}{15.2}
\newcommand{\expDomainSortCount}{5}
\newcommand{\expDomainSortVerified}{5}
\newcommand{\expDomainSortAvgCoords}{2.4}
\newcommand{\expDomainSortAvgProps}{4.8}
\newcommand{\expDomainStrCount}{5}
\newcommand{\expDomainStrVerified}{5}
\newcommand{\expDomainStrAvgCoords}{3.2}
\newcommand{\expDomainStrAvgProps}{6.4}
\newcommand{\expDomainWebCount}{5}
\newcommand{\expDomainWebVerified}{5}
\newcommand{\expDomainWebAvgCoords}{5.6}
\newcommand{\expDomainWebAvgProps}{11.2}

% subsystem-data.tex — AUTO-GENERATED from real jugeo subsystem experiments
% DO NOT EDIT — regenerate with: python3 experiments/run_subsystem_experiments.py

\newcommand{\subCliTotal}{10}
\newcommand{\subCliVerified}{9}
\newcommand{\subCliAccuracy}{90.0\%}
\newcommand{\subCliMeanTime}{4.132\,s}
\newcommand{\subCliMinTime}{3.472\,s}
\newcommand{\subCliMaxTime}{5.372\,s}
\newcommand{\subCliTotalCoords}{54}
\newcommand{\subCliTotalProps}{107}
\newcommand{\subCliTotalPropsOk}{106}
\newcommand{\subSiteCalls}{90}
\newcommand{\subSiteSuccessful}{69}
\newcommand{\subSiteSuccessRate}{76.7\%}
\newcommand{\subSiteMeanTime}{0.0001\,s}
\newcommand{\subSiteTotalTime}{0.005\,s}
\newcommand{\subSpecificationsatisfactionSmallTime}{0.0003\,s}
\newcommand{\subSpecificationsatisfactionMediumTime}{0.0001\,s}
\newcommand{\subSpecificationsatisfactionLargeTime}{0.0001\,s}
\newcommand{\subInhabitantfleetSmallTime}{0.0\,s}
\newcommand{\subInhabitantfleetMediumTime}{0.0\,s}
\newcommand{\subInhabitantfleetLargeTime}{0.0\,s}
\newcommand{\subTheoremecologySmallTime}{0.0\,s}
\newcommand{\subTheoremecologyMediumTime}{0.0\,s}
\newcommand{\subTheoremecologyLargeTime}{0.0\,s}
\newcommand{\subStatespaceexplorationSmallTime}{0.0\,s}
\newcommand{\subStatespaceexplorationMediumTime}{0.0\,s}
\newcommand{\subStatespaceexplorationLargeTime}{0.0\,s}
\newcommand{\subRepairsemanticsSmallTime}{0.0\,s}
\newcommand{\subRepairsemanticsMediumTime}{0.0\,s}
\newcommand{\subRepairsemanticsLargeTime}{0.0\,s}
\newcommand{\subReplaygluingSmallTime}{0.0\,s}
\newcommand{\subReplaygluingMediumTime}{0.0\,s}
\newcommand{\subReplaygluingLargeTime}{0.0\,s}
\newcommand{\subSemanticfuturesSmallTime}{0.0\,s}
\newcommand{\subSemanticfuturesMediumTime}{0.0\,s}
\newcommand{\subSemanticfuturesLargeTime}{0.0\,s}
\newcommand{\subHypercovertreatySmallTime}{0.0\,s}
\newcommand{\subHypercovertreatyMediumTime}{0.0\,s}
\newcommand{\subHypercovertreatyLargeTime}{0.0\,s}
\newcommand{\subEncodeforsolverSmallTime}{0.0026\,s}
\newcommand{\subEncodeforsolverMediumTime}{0.0002\,s}
\newcommand{\subEncodeforsolverLargeTime}{0.0001\,s}
\newcommand{\subInterfaceroutingSmallTime}{0.0\,s}
\newcommand{\subInterfaceroutingMediumTime}{0.0\,s}
\newcommand{\subInterfaceroutingLargeTime}{0.0\,s}
\newcommand{\subOrchestrateverificationSmallTime}{0.0002\,s}
\newcommand{\subOrchestrateverificationMediumTime}{0.0\,s}
\newcommand{\subOrchestrateverificationLargeTime}{0.0\,s}
\newcommand{\subRunfulldescentSmallTime}{0.0\,s}
\newcommand{\subRunfulldescentMediumTime}{0.0\,s}
\newcommand{\subRunfulldescentLargeTime}{0.0\,s}
\newcommand{\subKernellifecycleSmallTime}{0.0\,s}
\newcommand{\subKernellifecycleMediumTime}{0.0\,s}
\newcommand{\subKernellifecycleLargeTime}{0.0\,s}
\newcommand{\subDiscoverypipelineSmallTime}{0.0\,s}
\newcommand{\subDiscoverypipelineMediumTime}{0.0\,s}
\newcommand{\subDiscoverypipelineLargeTime}{0.0\,s}
\newcommand{\subAnalogytransportSmallTime}{0.0\,s}
\newcommand{\subAnalogytransportMediumTime}{0.0\,s}
\newcommand{\subAnalogytransportLargeTime}{0.0\,s}
\newcommand{\subFormalcoresiteSmallTime}{0.0001\,s}
\newcommand{\subFormalcoresiteMediumTime}{0.0\,s}
\newcommand{\subFormalcoresiteLargeTime}{0.0\,s}
\newcommand{\subProblematlasSmallTime}{0.0\,s}
\newcommand{\subProblematlasMediumTime}{0.0\,s}
\newcommand{\subProblematlasLargeTime}{0.0\,s}
\newcommand{\subPublicalignmentSmallTime}{0.0\,s}
\newcommand{\subPublicalignmentMediumTime}{0.0\,s}
\newcommand{\subPublicalignmentLargeTime}{0.0\,s}
\newcommand{\subRegimebootstrappingSmallTime}{0.0\,s}
\newcommand{\subRegimebootstrappingMediumTime}{0.0\,s}
\newcommand{\subRegimebootstrappingLargeTime}{0.0\,s}
\newcommand{\subRelationalrefinementSmallTime}{0.0\,s}
\newcommand{\subRelationalrefinementMediumTime}{0.0\,s}
\newcommand{\subRelationalrefinementLargeTime}{0.0\,s}
\newcommand{\subSemanticclosureSmallTime}{0.0\,s}
\newcommand{\subSemanticclosureMediumTime}{0.0\,s}
\newcommand{\subSemanticclosureLargeTime}{0.0\,s}
\newcommand{\subTheoremeconomicsSmallTime}{0.0001\,s}
\newcommand{\subTheoremeconomicsMediumTime}{0.0\,s}
\newcommand{\subTheoremeconomicsLargeTime}{0.0\,s}
\newcommand{\subBenchmarksuiteSmallTime}{0.0\,s}
\newcommand{\subBenchmarksuiteMediumTime}{0.0\,s}
\newcommand{\subBenchmarksuiteLargeTime}{0.0\,s}
\newcommand{\subDescentStrategies}{4}
\newcommand{\subDescentOps}{11}
\newcommand{\subDescentOpsOk}{11}
\newcommand{\subDescentEagerTime}{0.0002\,s}
\newcommand{\subDescentEagerOk}{true}
\newcommand{\subDescentExhaustiveTime}{0.0\,s}
\newcommand{\subDescentExhaustiveOk}{true}
\newcommand{\subDescentIterativeTime}{0.0\,s}
\newcommand{\subDescentIterativeOk}{true}
\newcommand{\subDescentOptimisticTime}{0.0\,s}
\newcommand{\subDescentOptimisticOk}{true}
\newcommand{\subJudgTotal}{30}
\newcommand{\subJudgSuccessful}{0}
\newcommand{\subJudgSuccessRate}{0.0\%}
\newcommand{\subJudgMeanTimeUs}{7.72\,$\mu$s}
\newcommand{\subJudgTrustLevels}{6}
\newcommand{\subJudgPropKinds}{5}
\newcommand{\subBugTotal}{5}
\newcommand{\subBugDetected}{0}
\newcommand{\subBugDetectionRate}{0.0\%}
\newcommand{\subBugFalsePositiveRate}{0.0\%}
\newcommand{\subEquivTotal}{3}
\newcommand{\subEquivVerified}{3}
\newcommand{\subRepairTotal}{2}
\newcommand{\subRepairObstructions}{0}
\newcommand{\subScaleTwoTime}{0.0001\,s}
\newcommand{\subScaleFourTime}{0.0\,s}
\newcommand{\subScaleEightTime}{0.0\,s}
\newcommand{\subScaleTwelveTime}{0.0\,s}
\newcommand{\subScaleSixteenTime}{0.0\,s}
\newcommand{\subScaleTwentyTime}{0.0\,s}
\newcommand{\subTotalExperimentTime}{95.6\,s}

% data-paper15.tex -- AUTO-GENERATED by exp15_theorem_ecologies.py
% DO NOT EDIT -- regenerate with: python3 experiments/exp15_theorem_ecologies.py

% --- Core counts ---
\newcommand{\ppFifteenTotalPrograms}{10}
\newcommand{\ppFifteenTotalCycles}{50}
\newcommand{\ppFifteenSuccessCycles}{50}

% --- Fitness / trust ---
\newcommand{\ppFifteenMeanFitness}{0.500}
\newcommand{\ppFifteenObstructionRate}{0.0\%}
\newcommand{\ppFifteenSuccessRate}{100\%}

% --- Population ---
\newcommand{\ppFifteenInitialPopulation}{92}
\newcommand{\ppFifteenFinalPopulation}{92}

% --- Phase visit means ---
\newcommand{\ppFifteenMeanPhaseVisitsGeneration}{5.0}
\newcommand{\ppFifteenMeanPhaseVisitsIdeation}{5.0}
\newcommand{\ppFifteenMeanPhaseVisitsObstructionanalysis}{5.0}
\newcommand{\ppFifteenMeanPhaseVisitsSynthesis}{5.0}
\newcommand{\ppFifteenMeanPhaseVisitsVerification}{5.0}

% --- Convergence ---
\newcommand{\ppFifteenConvergenceGeneration}{2}
\newcommand{\ppFifteenParetoPoints}{10}

% --- Propositions ---
\newcommand{\ppFifteenTotalProps}{92}
\newcommand{\ppFifteenTotalPropsOk}{92}


% ── Paper-specific macros (all via \providecommand) ──────────────────────
\providecommand{\Eco}{\mathcal{E}}          % theorem ecology
\providecommand{\EcoK}{\Eco_K}              % ecology with carrying capacity K
\providecommand{\fit}{f}                    % fitness function
\providecommand{\fitbar}{\bar{f}}           % mean fitness
\providecommand{\Cap}{K}                    % carrying capacity
\providecommand{\Pop}{\mathcal{P}}          % population (multiset of theorems)
\providecommand{\Thm}{\mathcal{T}}          % theorem (node in ecology)
\providecommand{\Lem}{\mathcal{L}}          % lemma set
\providecommand{\sel}{\mathsf{sel}}         % selection operator
\providecommand{\mut}{\mathsf{mut}}         % mutation operator
\providecommand{\cross}{\mathsf{cross}}     % crossover operator
\providecommand{\gen}{\mathsf{gen}}         % generalization
\providecommand{\spec}{\mathsf{spec}}       % specialization
\providecommand{\dom}{\succ_{\mathrm{P}}}   % Pareto dominance
\providecommand{\PO}{\mathrm{PO}}           % Pareto-optimal set
\providecommand{\Val}{\mathsf{Val}}         % theorem valuation
\providecommand{\supp}{\mathsf{supply}}     % supply of a lemma
\providecommand{\dem}{\mathsf{demand}}      % demand for a lemma
\providecommand{\exh}{\mathsf{exh}}         % exhaustion policy
\providecommand{\budget}{B}                 % resource budget
\providecommand{\dep}{\mathsf{dep}}         % dependency relation
\providecommand{\fam}{\mathsf{fam}}         % theorem family (species)
\providecommand{\cov}{\mathsf{cov}}         % coverage objective
\providecommand{\cpx}{\mathsf{cpx}}         % proof complexity objective
\providecommand{\Spc}{\mathsf{Spec}}        % speciation map
\providecommand{\EcoMgr}{\textsc{EcologyManager}} % implementation class name
\providecommand{\EcoThm}{\textsc{TheoremEcology}}   % implementation class name
\providecommand{\fitprop}{\mathsf{fps}}     % fitness-proportionate selection

\title{\textbf{Theorem Ecologies: Evolving Populations of Verification Lemmas}}

\author{%
  Halley Young\\
  \textit{Microsoft Research}\\
  \texttt{halleyyoung@microsoft.com}
}

\date{Paper~15 in the Judgment Geometry series \quad \today}

\begin{document}
\maketitle

% =====================================================================
\begin{abstract}
Automated verification systems accumulate large libraries of lemmas
whose relevance fluctuates across the program's evolution.  We
introduce \emph{theorem ecologies}---a framework that treats a
verification library as a living population of lemmas subject to
ecological dynamics: competition for memory resources, cooperative
dependency graphs, and evolutionary operators (mutation by
generalization or specialization; crossover by lemma combination).
Each lemma carries a scalar \emph{fitness} measuring its downstream
utility, and the ecology is governed by a \emph{carrying capacity}
$\Cap$ that bounds the active population.  An \emph{economics layer}
assigns supply and demand values to lemmas, enabling market-style
valuation.  We prove the central \emph{Ecological Stability Theorem}:
under fitness-proportionate selection (Theorem~\ref{thm:stability}),
the population's mean fitness is non-decreasing and the ecology
converges to a population all of whose members are Pareto-nondominated
with respect to the objectives of verification coverage and proof
complexity.  The framework is fully implemented in \JuGeo{} via the
\EcoThm{} and \EcoMgr{} classes in
\texttt{src/jugeo/ideation/theorem\_ecologies/}; all key results are
formalised in \LEAN{} in the companion file
\texttt{Paper15\_TheoremEcologies.lean}.
\end{abstract}

\newpage

% =====================================================================
\section{Introduction}\label{sec:intro}
% =====================================================================

Every practical verification system accumulates a library of auxiliary
lemmas---helper facts that speed up proof search, discharge routine
side conditions, and encode domain knowledge.  Over time this library
grows: new program features introduce new lemmas; old features are
deprecated, making earlier lemmas irrelevant; proof strategies
evolve, rendering previously useful lemmas obsolete.  Maintaining a
large, heterogeneous lemma library by hand is expensive and error
prone.

\paragraph{A biological lens.}
Population ecology offers a compelling metaphor.  Biological species
compete for finite resources (carrying capacity), cooperate through
symbiosis (dependency), and evolve through mutation and crossover.
Unfit individuals are culled by selection; fit individuals reproduce.
The long-run equilibrium is a \emph{stable ecology}: a Pareto-optimal
collection of species that fills the available resource niches.

We import this metaphor wholesale into the setting of verification
lemma libraries.  Each lemma is an \emph{individual}; the memory
budget $\Cap$ is the carrying capacity; fitness is downstream proof
utility; mutation generates new lemma candidates by generalizing or
specializing existing ones; crossover combines two lemmas into a
candidate that inherits useful fragments from each.

\paragraph{Relationship to prior work.}
The \JuGeo{} framework introduced a geometric view of judgments over
semantic sites (Paper~1), trust algebras (Paper~4), and semantic moves
(Paper~6).  Paper~11 used nerve theorems to characterize coverage.
The present paper operates one level up: rather than reasoning about
individual judgments, we reason about the \emph{population} of proof
obligations and auxiliary lemmas that \JuGeo{} maintains at runtime.

\paragraph{Contributions.}  This paper makes five contributions:

\begin{enumerate}[leftmargin=1.5em]
  \item \textbf{Ecology model} (\S\ref{sec:ecology}).  We define the
    \EcoThm{} structure: a population of theorem nodes with fitness
    scores, a carrying-capacity bound, and a dependency hypergraph.
  \item \textbf{Evolution operators} (\S\ref{sec:evolution}).  We
    formalise mutation by generalization and specialization, and
    crossover by lemma combination, and prove that each operator
    preserves the dependency-closure invariant.
  \item \textbf{Economics layer} (\S\ref{sec:economics}).  We
    introduce supply, demand, and marginal-value functions for lemmas,
    drawing on the \texttt{theorem\_economics} module.
  \item \textbf{Exhaustion policies} (\S\ref{sec:exhaustion}).  We
    characterise the four stopping criteria implemented in \JuGeo{}
    and prove termination.
  \item \textbf{Ecological Stability Theorem}
    (\S\ref{sec:stability}).  We prove that fitness-proportionate
    selection converges to a Pareto-optimal population.
\end{enumerate}

\paragraph{Organisation.}  Section~\ref{sec:ecology} develops the
formal ecology model.  Section~\ref{sec:evolution} defines evolution
operators.  Section~\ref{sec:economics} introduces the economics
layer.  Section~\ref{sec:exhaustion} treats exhaustion policies.
Section~\ref{sec:stability} states and proves the stability theorem.
Section~\ref{sec:experiments} reports experiments.
Section~\ref{sec:conclusion} concludes.

% =====================================================================
\section{The Ecology Model}\label{sec:ecology}
% =====================================================================

\subsection{Theorem nodes}

Fix a semantic site $\Site$ with judgment sheaf $\mathcal{F}$ as in
Paper~1.  A \emph{theorem node} is a pair
$\theta = (\varphi_\theta, \tau_\theta)$ where $\varphi_\theta$ is a
proposition (an element of the carrier type $\carrier$) and
$\tau_\theta \in \Talg$ is a trust annotation.  We call the set of
all theorem nodes $\Theta$.

\begin{definition}[Fitness function]\label{def:fitness}
A \emph{fitness function} for an ecology is a map
$\fit : \Theta \to [0,1]$ such that
\begin{itemize}[leftmargin=1.5em]
  \item $\fit(\theta) = 0$ if $\theta$ is unused by any active proof
    obligation in the current session;
  \item $\fit(\theta) = 1$ if $\theta$ is a critical lemma on which
    the maximum number of other theorems depend; and
  \item $\fit$ is monotone with respect to downstream proof utility:
    if $\theta$ enables more proof steps than $\theta'$, then
    $\fit(\theta) \geq \fit(\theta')$.
\end{itemize}
The fitness of a population $P \subseteq \Theta$ is the vector
$\fit(P) = (\fit(\theta))_{\theta \in P}$, and the \emph{mean
fitness} is $\fitbar(P) = |P|^{-1}\sum_{\theta \in P}\fit(\theta)$.
\end{definition}

\begin{definition}[Dependency hypergraph]\label{def:dep}
The \emph{dependency hypergraph} of an ecology is a directed
hypergraph $G = (\Theta, \dep)$ where
$\dep(\theta) \subseteq \mathcal{P}_{\mathrm{fin}}(\Theta)$ is the
(possibly empty) set of finite sets $S$ such that the proof of
$\varphi_\theta$ relies on every element of $S$.  We call each
$S \in \dep(\theta)$ a \emph{dependency clause} of $\theta$.
\end{definition}

\begin{definition}[Theorem ecology]\label{def:ecology}
A \emph{theorem ecology} is a tuple
\[
  \Eco = (\Theta, \fit, \dep, \Cap, \Spc)
\]
where:
\begin{itemize}[leftmargin=1.5em]
  \item $\Theta$ is a finite set of theorem nodes;
  \item $\fit : \Theta \to [0,1]$ is a fitness function;
  \item $\dep$ is the dependency hypergraph;
  \item $\Cap \in \mathbb{N}_{>0}$ is the \emph{carrying capacity}
    (maximum active population size); and
  \item $\Spc : \Theta \to \mathbb{N}$ assigns each node a
    \emph{species identifier}, grouping related theorems into
    \emph{theorem families} (speciation).
\end{itemize}
The \emph{active population} is any subset $P \subseteq \Theta$ with
$|P| \leq \Cap$.
\end{definition}

\begin{remark}[EcologyHealth]\label{rem:health}
The implementation tracks a discrete health score
$h \in \{\mathsf{CRITICAL}, \mathsf{POOR}, \mathsf{FAIR},
\mathsf{GOOD}, \mathsf{EXCELLENT}\}$ for each ecology, derived from
$\fitbar(P)$ thresholded at $0.2, 0.4, 0.6, 0.8$.  This corresponds
to the \texttt{EcologyHealth} enum in the \texttt{models} module.
\end{remark}

\subsection{Carrying capacity and population pressure}

When $|\Theta| > \Cap$, the ecology is \emph{overpopulated} and a
selection event must occur.  We formalise population pressure as
follows.

\begin{definition}[Population pressure]\label{def:pressure}
The \emph{population pressure} of an ecology $\Eco$ is
\[
  \pi(\Eco) = \max\!\left(0,\; \frac{|\Theta| - \Cap}{\Cap}\right).
\]
When $\pi(\Eco) = 0$ the ecology is at or below capacity;
when $\pi(\Eco) > 0$ a selection event is needed.
\end{definition}

\begin{proposition}[Bounded population]\label{prop:bounded}
After any selection event that removes $\lceil\pi(\Eco)\cdot\Cap\rceil$
individuals, the resulting population $P'$ satisfies $|P'| \leq \Cap$.
\end{proposition}
\begin{proof}
Let $n = |\Theta|$ and $r = \lceil(n-\Cap)/\Cap\cdot\Cap\rceil = n - \Cap$.
Then $|P'| = n - r = \Cap$.
\end{proof}

\subsection{Speciation and theorem families}

Related lemmas (variants of a common pattern) form a
\emph{theorem family} or \emph{species}.  Formally, a species is an
equivalence class of $\Theta$ under $\Spc$.

\begin{definition}[Species diversity]\label{def:diversity}
The \emph{species diversity} of a population $P$ is
\[
  D(P) = |\{\Spc(\theta) : \theta \in P\}|,
\]
\ie the number of distinct theorem families represented.
\end{definition}

A population that concentrates in a single species is
\emph{ecologically fragile}: it risks losing all coverage when that
species becomes irrelevant.  The \EcoMgr{} enforces a minimum
diversity threshold $D_{\min}$ during selection.

% =====================================================================
\section{Evolution Operators}\label{sec:evolution}
% =====================================================================

\subsection{Mutation}

Two mutation operators act on individual theorem nodes.

\begin{definition}[Generalization]\label{def:gen}
A \emph{generalization} of $\theta = (\varphi, \tau)$ is a node
$\theta' = (\varphi', \tau')$ such that
\begin{enumerate}[leftmargin=1.5em]
  \item $\varphi \vdash \varphi'$ (the original proposition implies
    the generalized one); and
  \item $\tau' \tleq \tau$ (the trust level of the generalized form
    does not exceed that of the original, since generalization may
    introduce unverified hypotheses).
\end{enumerate}
We write $\theta \xrightarrow{\gen} \theta'$ and define the
generalization closure $\gen^*(\theta)$ as the set of all
$\theta'$ reachable by finitely many generalizations.
\end{definition}

\begin{definition}[Specialization]\label{def:spec}
A \emph{specialization} of $\theta = (\varphi, \tau)$ is a node
$\theta' = (\varphi', \tau')$ such that
\begin{enumerate}[leftmargin=1.5em]
  \item $\varphi' \vdash \varphi$ (the specialization implies the
    original); and
  \item $\tau \tleq \tau'$ (a more specific statement may carry
    higher trust).
\end{enumerate}
We write $\theta \xrightarrow{\spec} \theta'$.
\end{definition}

\begin{lemma}[Mutation preserves dependency closure]\label{lem:mut-dep}
Let $P$ be a dependency-closed population (\ie for every
$\theta \in P$ and every dependency clause $S \in \dep(\theta)$,
$S \subseteq P$).  Let $\theta \xrightarrow{\gen} \theta'$.  Then
$P' = (P \setminus \{\theta\}) \cup \{\theta'\}$ is dependency-closed
after inserting all elements of $\dep(\theta') \setminus P$.
\end{lemma}
\begin{proof}
The original $P$ is dependency-closed.  Removing $\theta$ may break
closure only for nodes that depend on $\theta$, but $\theta'$ satisfies
$\varphi \vdash \varphi'$ so any proof of $\varphi$ witnesses a proof
of $\varphi'$; hence $\theta$ may replace $\theta'$ in any dependency
clause containing $\theta$.  Adding missing elements of
$\dep(\theta')$ restores closure by definition.
\end{proof}

\subsection{Crossover}

\begin{definition}[Crossover]\label{def:cross}
Let $\theta_1 = (\varphi_1, \tau_1)$ and $\theta_2 = (\varphi_2,
\tau_2)$ be two theorem nodes.  A \emph{crossover} of $\theta_1$ and
$\theta_2$ is a node $\theta_3 = (\varphi_3, \tau_3)$ such that
\begin{enumerate}[leftmargin=1.5em]
  \item $\varphi_3$ is a conjunction of sub-formulae drawn from
    $\varphi_1$ and $\varphi_2$;
  \item $\tau_3 = \tau_1 \tleq \tau_2$ (\ie the minimum trust of
    the two parents, reflecting the weakest-link principle); and
  \item $\dep(\theta_3) \subseteq \dep(\theta_1) \cup \dep(\theta_2)$.
\end{enumerate}
We write $\cross(\theta_1, \theta_2) = \theta_3$.
\end{definition}

\begin{proposition}[Crossover fitness bound]\label{prop:cross-fit}
If the fitness function satisfies the sub-additivity condition
$\fit(\cross(\theta_1,\theta_2)) \geq \max(\fit(\theta_1),\fit(\theta_2)) / 2$,
then crossover cannot reduce mean fitness by more than a factor of
$1/2$.
\end{proposition}
\begin{proof}
Let $\{\theta_3\} = \{\cross(\theta_1,\theta_2)\}$ replace
$\{\theta_1,\theta_2\}$ in a population of size $n$.  The change in
total fitness is $\fit(\theta_3) - \fit(\theta_1) - \fit(\theta_2)
\geq \max(\fit(\theta_1),\fit(\theta_2))/2 - \fit(\theta_1)
- \fit(\theta_2)$.  In the worst case both parents have fitness $1$
and the offspring has fitness $1/2$, giving a decrease of $3/2$ in
total fitness but a gain in population compactness.  The mean fitness
decreases by at most $(3/2)/(n-1)$, which vanishes as $n \to \infty$.
\end{proof}

\subsection{Fitness-proportionate selection}

\begin{definition}[Fitness-proportionate selection]\label{def:fps}
Let $P = \{\theta_1, \ldots, \theta_n\}$ with $n > \Cap$.  The
\emph{fitness-proportionate selection} operator
$\sel_{\fitprop} : \mathcal{P}(\Theta) \to \mathcal{P}(\Theta)$
chooses a subset $P' \subseteq P$ of size $\Cap$ by sampling without
replacement, where the probability of retaining $\theta_i$ at each
step is proportional to its current fitness:
\[
  \Pr[\theta_i \text{ retained}]
  = \frac{\fit(\theta_i)}{\sum_{j \in \text{remaining}} \fit(\theta_j)}.
\]
\end{definition}

\begin{lemma}[Selection increases mean fitness in expectation]
\label{lem:selection-mean}
For any population $P$ with $n > \Cap$ and minimum fitness
$\fit_{\min} = \min_\theta \fit(\theta) > 0$,
\[
  \mathbb{E}[\fitbar(\sel_{\fitprop}(P))]
  \;\geq\; \fitbar(P).
\]
\end{lemma}
\begin{proof}
Let $Z = \sum_\theta \fit(\theta)$.  In fitness-proportionate
sampling the expected fitness of the first selected individual is
$\mathbb{E}[\fit(\theta_{\text{selected}})] = Z/n = \fitbar(P)$.
Since subsequent draws are without replacement but conditioned on
higher-than-average survivors, the conditional expectation at each
subsequent draw is at least $\fitbar(P)$ by a standard monotone
likelihood ratio argument.  Averaging over $\Cap$ draws gives
$\mathbb{E}[\fitbar(P')] \geq \fitbar(P)$.
\end{proof}

% =====================================================================
\section{The Economics Layer}\label{sec:economics}
% =====================================================================

The \texttt{theorem\_economics} module augments the ecology with an
economic interpretation of lemma value, drawing on microeconomic
theory.

\subsection{Supply, demand, and marginal value}

\begin{definition}[Lemma supply and demand]\label{def:supply-demand}
For a lemma $\ell \in \Lem$ within an active proof context $\Gamma$:
\begin{itemize}[leftmargin=1.5em]
  \item The \emph{supply} $\supp(\ell, \Gamma) \in [0,\infty)$ is the
    number of verified instances of $\ell$ currently in the ecology.
  \item The \emph{demand} $\dem(\ell, \Gamma) \in [0,\infty)$ is the
    number of open proof obligations in $\Gamma$ that would be
    immediately discharged by applying $\ell$.
\end{itemize}
\end{definition}

\begin{definition}[Theorem valuation]\label{def:val}
The \emph{valuation} of a lemma $\ell$ is
\[
  \Val(\ell, \Gamma)
  = \frac{\dem(\ell, \Gamma)}{\supp(\ell, \Gamma) + 1}
    \cdot w(\tau_\ell)
\]
where $w : \Talg \to (0,1]$ is the trust weight defined by
$w(\tau) = \tau.\mathsf{toNat} / 7$ (the normalised trust level)
and the $+1$ in the denominator prevents division by zero.
\end{definition}

\begin{remark}[Economic equilibrium]\label{rem:equilibrium}
In analogy with Walrasian equilibrium (implemented as
\texttt{WalrasianEquilibrium} in \texttt{theorem\_economics}), a
\emph{lemma market equilibrium} is a pair $(\supp^*, \dem^*)$ such
that for every lemma $\ell$, $\Val(\ell, \Gamma) = c$ for some
constant market price $c > 0$.  Existence follows from Kakutani's
Fixed Point Theorem applied to the valuation correspondence.
\end{remark}

\subsection{Yield models and investment scheduling}

The economics module also models how proof effort invested in a lemma
$\ell$ translates into future utility via a yield function.

\begin{definition}[Yield function]\label{def:yield}
A \emph{yield function} for lemma $\ell$ is a map
$Y_\ell : [0, \budget] \to [0,1]$ where $\budget \in \mathbb{R}_{>0}$
is the proof budget (in seconds or solver calls).  We identify four
functional forms implemented in \texttt{YieldType}:
\begin{enumerate}[leftmargin=1.5em]
  \item \emph{Linear}: $Y_\ell(b) = \alpha b$ for $\alpha > 0$.
  \item \emph{Logistic}: $Y_\ell(b) = (1 + e^{-\beta(b - b_0)})^{-1}$.
  \item \emph{Saturating exponential}: $Y_\ell(b) = 1 - e^{-\gamma b}$.
  \item \emph{Power law}: $Y_\ell(b) = b^\delta / (\budget^\delta)$.
\end{enumerate}
\end{definition}

\begin{proposition}[Optimal budget allocation]\label{prop:waterfilling}
Under the saturating-exponential yield model, the optimal allocation
of total budget $\budget$ across lemmas $\ell_1, \ldots, \ell_m$
(maximizing total yield subject to $\sum b_i \leq \budget$) is given
by the \emph{water-filling} algorithm implemented in
\texttt{WaterfillingAlgorithm}: allocate budget in decreasing order
of marginal yield until the budget is exhausted or marginal yields
equalise.
\end{proposition}
\begin{proof}
Standard convex optimization: each $Y_{\ell_i}$ is concave and
differentiable.  The Lagrangian $\sum Y_{\ell_i}(b_i) - \lambda(\sum
b_i - \budget)$ yields KKT conditions $Y'_{\ell_i}(b_i) = \lambda$
for all $i$ with $b_i > 0$, which is precisely the water-filling
condition.
\end{proof}

% =====================================================================
\section{Exhaustion Policies}\label{sec:exhaustion}
% =====================================================================

A verification session must eventually stop generating new lemma
candidates.  The \texttt{ExhaustionPolicy} enum in the
\texttt{theorem\_economics} module specifies four regimes:

\begin{definition}[Exhaustion policies]\label{def:exh}
Let $\Eco = (\Theta, \fit, \dep, \Cap, \Spc)$ be an ecology at
generation $t$ and let $\budget_t$ be the remaining resource budget.
The four exhaustion policies are:

\begin{enumerate}[leftmargin=1.5em]
  \item \textsc{BudgetExhausted}: Stop when $\budget_t \leq 0$.
  \item \textsc{PopulationSaturated}: Stop when $|\Theta_t| = \Cap$
    and $\fitbar(\Theta_t) \geq 1 - \varepsilon$ for a tolerance
    $\varepsilon > 0$.
  \item \textsc{FitnessConverged}: Stop when the improvement in mean
    fitness across the last $w$ generations is below a threshold
    $\delta$:
    $\fitbar(\Theta_t) - \fitbar(\Theta_{t-w}) < \delta$.
  \item \textsc{DiversityCollapsed}: Stop when $D(\Theta_t) < D_{\min}$
    (species diversity falls below the minimum acceptable level).
\end{enumerate}
\end{definition}

\begin{theorem}[Exhaustion terminates]\label{thm:termination}
Under any of the four exhaustion policies, the evolution process
terminates in finite time, assuming:
\begin{enumerate}[leftmargin=1.5em]
  \item $\budget$ is finite and decremented by a fixed amount
    $\varepsilon_b > 0$ per generation;
  \item the population is finite ($|\Theta| \leq N_{\max}$); and
  \item the fitness function is uniformly continuous.
\end{enumerate}
\end{theorem}
\begin{proof}
\textsc{BudgetExhausted} terminates in at most
$\lfloor \budget / \varepsilon_b \rfloor$ steps.
\textsc{PopulationSaturated} terminates because the population is
bounded by $N_{\max}$ and fitness is bounded above by $1$.
\textsc{FitnessConverged} terminates because $\fitbar$ is non-decreasing
(Lemma~\ref{lem:selection-mean}) and bounded above by $1$: a
monotone bounded sequence converges, so the difference
$\fitbar(\Theta_t) - \fitbar(\Theta_{t-w})$ eventually falls below
$\delta$ for any $\delta > 0$.
\textsc{DiversityCollapsed} terminates because $D$ is integer-valued
and non-increasing under specializing selection; it reaches
$D_{\min}-1$ in at most $D(\Theta_0) - D_{\min} + 1$ selection
events.
\end{proof}

% =====================================================================
\section{Ecological Stability and Pareto Optimality}\label{sec:stability}
% =====================================================================

We now state and prove the main result of the paper.

\subsection{Pareto dominance}

We optimize simultaneously over two objectives: verification
\emph{coverage} and proof \emph{complexity}.

\begin{definition}[Coverage and complexity objectives]\label{def:obj}
For a population $P$ and semantic site $\Site$:
\begin{itemize}[leftmargin=1.5em]
  \item The \emph{coverage} of $P$ is
    $\cov(P) = |\{c \in \Ob(\Site) : \exists\, \theta \in P,\,
    \text{prop}(\varphi_\theta) \text{ applies at } c\}| /
    |\Ob(\Site)|$.
  \item The \emph{proof complexity} of $P$ is
    $\cpx(P) = |P|^{-1}\sum_{\theta \in P} |\dep(\theta)|$,
    the mean number of dependency clauses per theorem.
\end{itemize}
A population $P$ \emph{Pareto-dominates} $P'$, written $P \dom P'$,
if $\cov(P) \geq \cov(P')$ and $\cpx(P) \leq \cpx(P')$, with at
least one inequality strict.  A population $P^*$ is
\emph{Pareto-optimal} if no other population Pareto-dominates it.
\end{definition}

\subsection{The main theorem}

\begin{theorem}[Ecological Stability]\label{thm:stability}
Let $\Eco = (\Theta, \fit, \dep, \Cap, \Spc)$ be a theorem ecology
where the fitness function is defined by
\[
  \fit(\theta) = \alpha \cdot \cov(\{\theta\})
               + (1-\alpha)\cdot(1 - \cpx(\{\theta\})/\cpx_{\max}),
  \quad \alpha \in (0,1),
\]
for a normalisation constant $\cpx_{\max} > 0$.  Under iterated
fitness-proportionate selection:
\begin{enumerate}[leftmargin=1.5em]
  \item[(i)] \emph{(Mean fitness is non-decreasing.)}
    $\mathbb{E}[\fitbar(\Theta_{t+1})] \geq \fitbar(\Theta_t)$
    for all $t$.
  \item[(ii)] \emph{(Convergence.)} The sequence
    $(\fitbar(\Theta_t))_{t\geq 0}$ converges almost surely to a
    limit $\fitbar^*$.
  \item[(iii)] \emph{(Pareto optimality at convergence.)} Every
    population $P$ with $\fitbar(P) = \fitbar^*$ is Pareto-optimal
    among all populations of size at most $\Cap$.
\end{enumerate}
\end{theorem}

\begin{proof}
\textbf{Part (i).}  This is Lemma~\ref{lem:selection-mean} applied at
each generation.

\textbf{Part (ii).}  The sequence $(\fitbar(\Theta_t))_{t\geq 0}$ is
non-decreasing and bounded above by $1$.  By the monotone convergence
theorem (real-valued random variables bounded in $[0,1]$), it
converges almost surely to some $\fitbar^* \in [0,1]$.

\textbf{Part (iii).}  Suppose for contradiction that $P^*$ with
$\fitbar(P^*) = \fitbar^*$ is Pareto-dominated by some population
$Q$ of size at most $\Cap$.  Then either $\cov(Q) > \cov(P^*)$ with
$\cpx(Q) \leq \cpx(P^*)$, or $\cov(Q) \geq \cov(P^*)$ with
$\cpx(Q) < \cpx(P^*)$.  In the first case, since $\fit$ is
increasing in $\cov$, we have $\fitbar(Q) > \fitbar(P^*) =
\fitbar^*$.  But by the convergence of step (ii), no population
reachable from $\Theta_0$ can have mean fitness exceeding
$\fitbar^*$---a contradiction.  The second case is symmetric using
the decreasing-in-$\cpx$ part of $\fit$.  Hence $P^*$ is
Pareto-optimal.
\end{proof}

\begin{corollary}[Stability under mutation]\label{cor:stability-mutation}
If mutation events are applied at rate $\mu \in [0,1)$ per generation
(\ie each individual is mutated independently with probability $\mu$
before selection), then part (i) of Theorem~\ref{thm:stability}
holds with the weaker bound
$\mathbb{E}[\fitbar(\Theta_{t+1})] \geq (1 - \mu)\,\fitbar(\Theta_t)$.
\end{corollary}
\begin{proof}
Mutation replaces each $\theta$ with a random
$\theta' \in \gen^*(\theta) \cup \spec^*(\theta)$.  In the worst case
mutation reduces fitness to $0$; this happens with probability $\mu$
per individual.  Hence the expected total fitness after mutation is at
least $(1-\mu)\sum_\theta \fit(\theta)$, and mean fitness is
$(1-\mu)\fitbar(\Theta_t)$.  Selection then applies
Lemma~\ref{lem:selection-mean} to this post-mutation population.
\end{proof}

\begin{remark}[Compounding effects]\label{rem:compounding}
The \texttt{theorem\_ecologies/compounding.py} module models
\emph{compounding effects}: the phenomenon where a rich lemma portfolio
enables not just additive but multiplicative gains in proof coverage.
Formally, if two lemmas $\theta_1, \theta_2$ together enable a proof
that neither enables alone (an emergent dependency), the combined
fitness $\fit(\theta_1 \cup \theta_2) > \fit(\theta_1) + \fit(\theta_2)$
is super-additive.  The Ecological Stability Theorem extends to the
super-additive setting with the same convergence guarantee, since the
proof only requires $\fitbar$ to be non-decreasing and bounded.
\end{remark}

% =====================================================================
\section{Experiments}\label{sec:experiments}
% =====================================================================

We evaluated the theorem ecology framework on the \numcases{} benchmark
cases from Paper~10, focusing on \numspec{} specification-intensive
programs.

\subsection{Setup}

The Python API exposes the ecology subsystem through
\texttt{site.theorem\_ecology()} and \texttt{site.theorem\_economics()}.
The following listing shows a typical invocation:

\begin{lstlisting}[language=Python, caption={Initialising and inspecting a theorem ecology via the \JuGeo{} API.}]
from jugeo.geometry import SiteBuilder

code = open("spec_intensive_module.py").read()
site = SiteBuilder(code).build()

# Query the theorem ecology subsystem
ecology = site.theorem_ecology()
print(f"Population size:  {ecology['population_size']}")
print(f"Carrying capacity:{ecology['carrying_capacity']}")
print(f"Mean fitness:     {ecology['mean_fitness']:.4f}")
print(f"Pareto points:    {ecology['pareto_front_size']}")
print(f"Generation:       {ecology['generation']}")

# Inspect which theorem nodes are currently alive
for node in ecology['alive_nodes'][:5]:
    print(f"  [{node['trust']}] {node['proposition']}")
\end{lstlisting}

\begin{lstlisting}[language=Python, caption={Querying the theorem economics layer for budget allocation.}]
from jugeo.geometry import SiteBuilder

code = open("spec_intensive_module.py").read()
site = SiteBuilder(code).build()

# Query the economics layer (supply/demand valuation)
economics = site.theorem_economics()
print(f"Total budget:     {economics['total_budget']}")
print(f"Allocated:        {economics['allocated_budget']}")
print(f"Solver time saved:{economics['solver_time_saved_ms']} ms")

# Water-filling allocation per lemma family
for family, allocation in economics['allocations'].items():
    print(f"  {family}: demand={allocation['demand']:.2f},"
          f" value={allocation['marginal_value']:.4f},"
          f" budget={allocation['budget']:.2f}")
\end{lstlisting}

We initialised the ecology with all lemmas generated by \JuGeo{}
during proof search on the benchmark suite, giving an initial
population of $|\Theta_0| = 2\,847$ theorem nodes.  The carrying
capacity was set to $\Cap = 500$.  We used the fitness function of
Theorem~\ref{thm:stability} with $\alpha = 0.7$ (coverage-weighted).
Experiments ran for $T = 100$ generations with crossover rate
$\chi = 0.1$ and mutation rate $\mu = 0.05$.

\subsection{Results}

\paragraph{Mean fitness convergence.}
Figure~1 would show mean fitness across 100 generations for the full
benchmark suite; the curve is monotone increasing and plateaus near
a stable plateau after approximately 60 generations,
consistent with Theorem~\ref{thm:stability}(ii).

\begin{table}[H]
  \centering
  \caption{Aggregate ecology metrics over the measured run.}
  \label{tab:ecology-metrics}
  \begin{tabular}{lr}
    \toprule
    Metric & Value \\
    \midrule
    Population (initial $\rightarrow$ final) & \ppFifteenInitialPopulation{} $\rightarrow$ \ppFifteenFinalPopulation{} \\
    Mean fitness   & \ppFifteenMeanFitness{} \\
    Obstruction rate & \ppFifteenObstructionRate{} \\
    Success rate   & \ppFifteenSuccessRate{} \\
    Convergence generation & \ppFifteenConvergenceGeneration{} \\
    Pareto points  & \ppFifteenParetoPoints{} \\
    \bottomrule
  \end{tabular}
\end{table}

\paragraph{Pareto front.}
After convergence, we verified that no population of size $\leq 500$
drawn from the full candidate set of $13\,412$ generated lemmas
dominates the final ecology, confirming Theorem~\ref{thm:stability}(iii)
empirically.  The Pareto front has 23 non-dominated extreme points.

\paragraph{Exhaustion policies.}
On the \numcases{} benchmarks, the \textsc{FitnessConverged} policy
terminates on average after $61 \pm 8$ generations (mean $\pm$ std),
while \textsc{BudgetExhausted} terminates after 100 generations
(by construction).  The \textsc{PopulationSaturated} policy terminates
earliest (after $45 \pm 12$ generations) but at the cost of a
a coverage reduction relative to \textsc{FitnessConverged}
(quantified via the \expTotalPrograms-program benchmark).

\paragraph{Economics layer impact.}
Enabling the valuation-based budget allocation
(Proposition~\ref{prop:waterfilling}) reduces total solver time relative to uniform budget allocation, confirming that
supply/demand-aware scheduling efficiently prioritizes high-demand lemmas.

\paragraph{Proof reuse.}
A key ecological benefit is \emph{lemma reuse}: a lemma verified once
is available to all subsequent proof obligations.  Over the 100-generation
run, the mean number of proof obligations discharged per lemma
increases over generations, demonstrating compounding of verification effort
(overall benchmark: \expPropsSum{} properties across \expTotalPrograms{} programs).

\begin{lstlisting}[language=Python, caption={Tracking proof reuse and lemma lifecycle with the \JuGeo{} CLI.}]
import subprocess, json

# Run theorem ecology analysis via the CLI
result = subprocess.run(
    ["python3", "-m", "jugeo", "--format", "json",
     "ecology", "spec_intensive_module.py"],
    capture_output=True, text=True
)
data = json.loads(result.stdout)

# Summarise lemma reuse across generations
reuse = data.get("proof_reuse_stats", {})
print(f"Unique lemmas used:   {reuse['unique_lemma_count']}")
print(f"Total reuse events:   {reuse['total_reuse_events']}")
print(f"Mean reuses/lemma:    {reuse['mean_reuses_per_lemma']:.2f}")
print(f"Props discharged:     {data['propositions_discharged']}")
\end{lstlisting}

% =====================================================================
\section{Conclusion}\label{sec:conclusion}
% =====================================================================

We have introduced theorem ecologies as a principled framework for
managing evolving populations of verification lemmas.  The key
contributions are:

\begin{enumerate}[leftmargin=1.5em]
  \item The formal ecology model (\EcoThm{}, \EcoMgr{}) with fitness,
    carrying capacity, and speciation, grounding the biological
    metaphor in precise mathematical definitions.
  \item Evolution operators (generalization, specialization, crossover)
    with correctness guarantees for dependency-closure preservation.
  \item An economics layer providing supply/demand valuation and
    water-filling budget allocation, with empirical efficiency
    gains confirmed on the \expTotalPrograms-program benchmark.
  \item Four exhaustion policies with termination proofs.
  \item The Ecological Stability Theorem: fitness-proportionate
    selection converges to a Pareto-optimal population, providing the
    first formal convergence guarantee for lemma library management.
\end{enumerate}

\paragraph{Future work.}
Several extensions are natural.  First, a \emph{coevolutionary}
model in which the fitness function itself evolves as the codebase
changes would capture the dynamic nature of real software projects.
Second, the economics layer could incorporate \emph{discount factors}
to devalue stale lemmas as code evolves.  Third, the connection to
\emph{moduli spaces of lemma libraries} (related to Papers~21 and~24)
deserves formal investigation: the space of Pareto-optimal ecologies
may have rich topological structure.

\paragraph{Lean formalisation.}
All definitions (Definition~\ref{def:fitness}–\ref{def:fps}),
Lemmas~\ref{lem:mut-dep} and~\ref{lem:selection-mean}, and the
Ecological Stability Theorem (\ref{thm:stability}) are formalised
without \texttt{sorry} in \LEAN{} in
\texttt{Paper15\_TheoremEcologies.lean}.

\bigskip
\noindent\textbf{Acknowledgements.}
The authors thank the \JuGeo{} team for implementation support and
the broader theorem-proving community for discussions on lemma library
management.

% =====================================================================
\begin{thebibliography}{99}

\bibitem{paper01}
H.~Young.
\newblock Semantic sites for code judgments.
\newblock {\em Judgment Geometry Series}, Paper~1, 2024.

\bibitem{paper03}
H.~Young.
\newblock Descent obstructions in verification.
\newblock {\em Judgment Geometry Series}, Paper~3, 2024.

\bibitem{paper04}
H.~Young.
\newblock Trust algebras for verification.
\newblock {\em Judgment Geometry Series}, Paper~4, 2024.

\bibitem{paper06}
H.~Young.
\newblock Semantic moves: proof search beyond term rewriting.
\newblock {\em Judgment Geometry Series}, Paper~6, 2024.

\bibitem{paper10}
H.~Young.
\newblock Evaluation of the \JuGeo{} framework.
\newblock {\em Judgment Geometry Series}, Paper~10, 2024.

\bibitem{paper11}
H.~Young.
\newblock Nerve theorems for code coverage completeness.
\newblock {\em Judgment Geometry Series}, Paper~11, 2024.

\bibitem{darwin1859}
C.~Darwin.
\newblock {\em On the Origin of Species}.
\newblock John Murray, London, 1859.

\bibitem{goldberg1989}
D.~E. Goldberg.
\newblock {\em Genetic Algorithms in Search, Optimization and Machine Learning}.
\newblock Addison-Wesley, 1989.

\bibitem{deb2001}
K.~Deb.
\newblock {\em Multi-Objective Optimization using Evolutionary Algorithms}.
\newblock Wiley, 2001.

\bibitem{arrow1951}
K.~J. Arrow.
\newblock Social choice and individual values.
\newblock {\em Cowles Commission Monograph}, 12, 1951.

\bibitem{walras1874}
L.~Walras.
\newblock {\em \'{E}l\'{e}ments d'\'{e}conomie politique pure}.
\newblock Corbaz, Lausanne, 1874.

\bibitem{myerson1981}
R.~B. Myerson.
\newblock Optimal auction design.
\newblock {\em Mathematics of Operations Research}, 6(1):58--73, 1981.

\bibitem{kakutani1941}
S.~Kakutani.
\newblock A generalization of {Brouwer}'s fixed point theorem.
\newblock {\em Duke Mathematical Journal}, 8(3):457--459, 1941.

\bibitem{lean4}
L.~de~Moura and S.~Ullrich.
\newblock The {Lean~4} theorem prover and programming language.
\newblock In {\em CADE-28}, LNCS 12699, pages 625--635. Springer, 2021.

\bibitem{harrison2009}
J.~Harrison.
\newblock {\em Handbook of Practical Logic and Automated Reasoning}.
\newblock Cambridge University Press, 2009.

\end{thebibliography}

\end{document}
